

\documentclass[a4paper,11pt,twoside,notitlepage,reqno]{amsart}
% "reqno" makes equation tags appear on right (instead of left)

\usepackage{amsmath,amssymb,amsbsy,amsthm}
\usepackage[hmargin=1.4in,vmargin=1.4in]{geometry}
\usepackage{url}

\usepackage{tikz}
\usepackage{tikz-cd}    % Required for the commutative diagram of the self-correspondence
\usepackage{amsmath}    % Required for the math environments, \text{}, and matrix/diagram layouts
\usepackage{amssymb}    % Required for mathematical symbols like \mathbb, \setminus, etc.
\usepackage{amsthm}     % Required for defining and formatting the Theorem and Definition environments
\usetikzlibrary{matrix}
% hyperlinks
\usepackage[linktocpage]{hyperref}
%\usepackage[linktocpage]{hyperref}
\hypersetup{
    colorlinks=true,        % false: boxed links; true: colored links
    linkcolor=red,          % color of internal links (change box color with linkbordercolor)
    citecolor=red,         % color of links to bibliography
    filecolor=magenta,      % color of file links
    urlcolor=cyan           % color of external links
}

%Chapter/heading format
% \usepackage{titlesec}
% \titleformat{\section}[block]
%   {\filcenter}{\thesection .}{0.5em}{}
% \titleformat{\subsection}[runin]
%   {\normalfont}{\bfseries \thesubsection}{0.5em}{\bfseries}
% \titlespacing*{\subsection}{0pt}{22pt}{0.5em}

% This removes subsections from tableofcontents
\addtocontents{toc}{\protect\setcounter{tocdepth}{1}}

% Changing \mathcal{C}
\usepackage{eucal}

\usepackage{algorithm}
\usepackage{algpseudocode}
\newcommand{\red}{\color{red}}
\newcommand{\blue}{\color{blue}}

% some useful commands
\newcommand{\C}{ \mathbb{C} }
\newcommand{\R}{ \mathbb{R} }
\newcommand{\Q}{ \mathbb{Q} }
\newcommand{\Z}{ \mathbb{Z} }
\newcommand{\N}{ \mathbb{N} }
\newcommand{\F}{ \mathbb{F} }
\newcommand{\A}{ \mathbb{A} }
\DeclareMathOperator{\GL}{GL}
\newcommand{\ngr}{\trianglerighteq}
\newcommand{\ngl}{\trianglelefteq}
\newcommand{\ngrneq}{\triangleright}
\newcommand{\nglneq}{\triangleleft}
\newcommand{\cl}[1]{ \overline{#1} }

\newcommand{\aA}{ \mathcal{A} }
\newcommand{\bB}{ \mathcal{B} }
\newcommand{\cC}{ \mathcal{C} }
\newcommand{\dD}{ \mathcal{D} }
\newcommand{\fF}{ \mathcal{F} }

% theorem/proof environments, numbering, etc
%\newtheoremstyle{slstyle}% name of the style to be used
%  {}% measure of space to leave above the theorem. E.g.: 3pt
%  {}% measure of space to leave below the theorem. E.g.: 3pt
%  {\slshape}% name of font to use in the body of the theorem
%  {0pt}% measure of space to indent
%  {\bfseries}% name of head font
%  {. }% punctuation between head and body
%  {1pt}% space after theorem head; " " = normal interword space
%  {}% Manually specify head

%\swapnumbers

\theoremstyle{plain}
	\newtheorem{thm}{Theorem}
	\newtheorem{quest}{Question}
	\newtheorem{notn}[thm]{Notation}
		\numberwithin{thm}{section}
	\newtheorem{lemma}[thm]{Lemma}
	\newtheorem{prop}[thm]{Proposition}
	\newtheorem{cor}[thm]{Corollary}
	\newtheorem{conj}[thm]{Conjecture}

	\newtheorem*{thm*}{Theorem}
	\newtheorem*{lemma*}{Lemma}
	\newtheorem*{prop*}{Proposition}
	\newtheorem*{cor*}{Corollary}
	\newtheorem*{conj*}{Conjecture}

\theoremstyle{definition}
	\newtheorem{example}[thm]{Example}
	\newtheorem*{example*}{Example}
	\newtheorem{defn}[thm]{Definition}
	\newtheorem{remark}[thm]{Remark}
	%\newtheorem{remark}[thm]{Remark}

% % % % % % % % % % % %

\begin{document}


\title[]
{Filtered deformations of three-variable polynomial algebras}


\author{Jason Bell}
\author{Boris Li}
\address{University of Waterloo \\
Department of Pure Mathematics \\
Waterloo, Ontario \\
Canada  N2L 3G1}
\thanks{The authors were supported in part by NSERC grant RGPIN-2022-02951.}
\email{jpbell@uwaterloo.ca}
\email{boris.li@uwaterloo.ca}
\keywords{filtered deformations, positive characteristic, polynomial identities.}
\subjclass{16S80, 16S38, 13A35}
\begin{abstract}
    We classify all filtered deformations of polynomial rings in three variables.  We use our characterization to answer a special case of a question of Etingof by showing that if $F$ is an algebraically closed field of positive characteristic and $A$ is an $F$-algebra with a filtration by finite dimensional subspaces whose associated graded ring is a polynomial ring in three variables then $A$ satisfies a polynomial identity.  \end{abstract}
\maketitle

\section{Introduction}
In this paper, we look at ``quantized'' algebras that are constructed by deforming the multiplicative structure of a commutative polynomial ring. A general way of building such quantized algebras comes from studying algebras that are filtered and whose associated graded algebra with respect to this filtration is exactly this polynomial ring.  Given a field $K$ and a graded $K$-algebra $B$, a $\mathbb{Z}_+$-\emph{filtered deformation} (henceforth we shall drop the $\mathbb{Z}_+$) of $B$ is simply an algebra $A$ with a filtration 
\begin{equation}\label{eq:filt}
V_0\subseteq V_1 \subseteq V_2 \subseteq \cdots \subseteq \bigcup_i V_i=A
\end{equation} such that $B$ is the associated graded algebra of $A$.  An important class of filtered deformations of polynomial rings comes from enveloping algebras over a base field $K$ where the filtration is given by $V_n=V^n$, and where $V$ is the subspace spanned by the image of the Lie algebra and $1$.  In positive characteristic, enveloping algebras of finite-dimensional Lie algebras are known to be well-behaved.  In particular, they satisfy a polynomial identity.  One can think of this as saying that filtered deformations of polynomial algebras that are generated in degree one behave much like algebras that are finite over their centres.  Etingof \cite[Question 1.1]{Etingof} asked whether this phenomenon holds more generally for all filtered deformations of graded commutative algebras over positive characteristic algebraically closed fields and he proved that when a filtered deformation of a commutative algebra in positive characteristic satisfies a polynomial identity then the degree of the identity is a power of the characteristic of the field.

Understanding when filtered deformations in positive characteristic satisfy a polynomial identity is often useful when looking at filtered deformations of finitely generated commutative algebras in characteristic zero. The reason for this is that the algebras are finitely presented and hence have models over a base ring that is finitely generated as a $\mathbb{Z}$-algebra.  One can then reduce mod maximal ideals of the base ring to obtain filtered deformations of algebras in positive characteristic, and the fact that these have large centres in positive characteristic can sometimes be used to understand subtle phenomena in the characteristic zero setting. In some circumstances, for instance, this approach can be used to lift mod $p$ information to prove non-trivial results in characteristic zero; this philosophy inspired a related question of Cuadra, Etingof, and Walton \cite{CEW}.   

In general, Etingof's question has proved difficult to answer, although it has been answered in the affirmative in the case of iterated Hopf Ore extensions in positive characteristic \cite{BZ} and for algebras of Gelfand-Kirillov dimension at most two \cite{Bell}.  Etingof \cite{Etingof} notes that the case of filtered deformations of polynomial rings is already non-trivial and is of special interest. In this paper, we prove that Etingof's question has an affirmative answer in the case when the associated graded algebra is a polynomial ring in three variables.  
\begin{thm}
Let $K$ be an algebraically closed field of positive characteristic and let $A$ be a $K$-algebra with an exhaustive filtration by finite-dimensional $K$-vector subspaces, whose associated graded ring is a polynomial algebra over $K$ in three variables.  Then $A$ satisfies a polynomial identity.
\label{thm:main2}
\end{thm}

Our approach to proving Theorem \ref{thm:main2} is admittedly ham-fisted and relies on first classifying filtered deformations of three-variable polynomial rings, although we believe this classification is interesting in its own right.  Even armed with a classification, showing the algebras satisfy polynomial identities in positive characteristic is non-trivial for several of the families we obtain.  The reason for this is that in many cases, the families have models over characteristic zero fields that don't satisfy a polynomial identity and so the PI degrees of the algebras grow with the characteristic of the field; moreover, in many cases there are not obvious central elements one can work with. 

An important and natural class of filtered deformations of polynomial rings comes from \emph{iterated Ore extensions of unipotent type}.  We recall that a skew polynomial ring over a base $K$-algebra $R$ is $R[x]$ as a set but with multiplication ``deformed'' by extending the multiplication on $R$ to $R[x]$ via 
$x\cdot r = \sigma(r) x +\delta(r)$ where $\sigma$ is a $K$-linear automorphism of $R$ and $\delta$ is a $\sigma$-derivation; i.e., $\delta$ is $K$-linear and $\delta(ab)=\delta(a)b+\sigma(a)\delta(b)$ for $a,b\in R$.  We write $R[x;\sigma,\delta]$ to denote the resulting algebra. We say that an automorphism $\sigma$ of $R$ is \emph{unipotent} if the operator $\Delta = \sigma-{\rm Id}$ is locally nilpotent. If $R$ is finitely generated and $K$ has characteristic $p>0$ then this says $\sigma$ has order a power of $p$, since $\sigma^{p^n} = {\rm Id} +\Delta^{p^n}$ and for $n$ large enough, $\Delta^{p^n}$ annihilates a generating set of $R$.  We then say that a ring is an iterated Ore extension of $K$ if it is isomorphic to an algebra of the form $K[x_1][x_2;\sigma_2,\delta_2]\cdots [x_n;\sigma_n,\delta_n]$ and if, moreover, each $\sigma_i$ is unipotent, we say that it is of unipotent type.

A careful study yields the following lengthy classification.  We note that many of the families only exist in characteristics $2$ or $3$.
\begin{thm}\label{thm:main}
Let $K$ be a field and let $A$ be a filtered deformation of a three-variable polynomial ring over $K$.  Then $A$ is isomorphic to one of the following algebras:
\begin{enumerate}
\item An iterated Ore extension $K[x][y;\sigma_1,\delta_1][z;\sigma_2,\delta_2]$ of $K$ of unipotent type;
\item The algebra with generators $x,y,z$ and relations $$[x,y]=z, \qquad [x,z]=\gamma y, \qquad [y,z]=f(x)$$ for some $\gamma\in K$ and some $f(x)\in K[x]$.




\item The algebra with generators $x,y,z$ and relations
$$[x,y]=z, \quad [x,z]=\delta\beta(z-2xy)+\delta y^2+\beta x^3+\gamma_2 x^2+\gamma_1 x+\gamma_0,$$ 
$$[y,z]=u(x)-\Big(\gamma_1+(2\gamma_2-2\delta\beta^2)x+3\beta x^2\Big)y+\Big(\gamma_2-\delta\beta^2+3\beta x\Big)z.$$
 $\delta, \beta,\gamma_0,\gamma_1,\gamma_2\in K$, $u(x)\in K[x]$ of degree at most four; moreover, we can take either $\delta=0$; $\delta\neq 0, \beta=\delta^{-1}$; $\delta\neq 0, \beta=0$.
 
\item The algebra with generators $x,y,z$ and relations
$$[x,y]=z, [x,z]= z -2 xy +\zeta, [y,z]= y^2 +u(x),$$
where $\zeta\in K$, $u(x)\in K[x]$.

\item When the characteristic of $K$ is $3$, we get the family with generators
$x,y,z$ and relations
$$[x,y]=z,\quad [x,z]= \alpha z +\delta\gamma\, xy + \delta y^2 + \big(\gamma x^3 +\delta \gamma^2 x^2-\gamma (\alpha-\delta \gamma)^2 x+\mu\big),$$
and
$$[y,z]= \gamma(\delta\gamma -\alpha)\,z,$$
where $\alpha,\delta,\gamma,\mu\in K$.

%\item When the characteristic of $K$ is $3$, we get the family with generators $x,y,z$ and relations
%$$[x,y]=z,\quad [x,z]= z +\delta\gamma xy + \delta y^2 + (\gamma x^3 +\delta \gamma^2 x^2-\gamma (1-\delta \gamma)^2 x+\mu),$$ and $$[y,z]= (\gamma^2\delta -\gamma)z,$$ where $\delta, \gamma,\mu\in K$.
%\item When the characteristic of $K$ is $2$, we get the family with generators $x,y,z$ and relations
%$$[x,y]=z, \quad [x,z]= \beta y +  \gamma_2 x^2 + \gamma_1 x + \gamma_0 +\delta y^2,\quad [y,z]=u(x)+\gamma_1y +\gamma_2 z,$$ where $\beta,\gamma_2,\gamma_1,\gamma_0,\delta \in K$ and $u(x)\in K[x]$ of degree at most four.
%\%item When the characteristic of $K$ is $2$, we get the family with generators $x,y,z$ and relations $$[x,y]=z, [x,z]=z + \beta y + (\delta^{-1} x^3+ \gamma_2 x^2 + \gamma_1 x + \gamma_0 ) +\delta y^2,$$
%and
%$$[y,z] = u(x) + (\delta^{-1} x^2+\gamma_1+\beta \delta^{-1})y + (\delta^{-1} x+\gamma_2+\delta^{-1})z,$$
%where $u(x)\in K[x]$ has degree at most $4$, $\delta, \beta, \gamma_2,\gamma_1,\gamma_0\in K$ where $\delta\neq 0$.

 \item When the characteristic of $K$ is $2$, we get the family with generators $x,y,z$ and relations
$$[x,y]=z, \quad [x,z]= \alpha z + \beta y + \big(\alpha\delta^{-1} x^3+\gamma_2 x^2 + \gamma_1 x + \gamma_0\big) +\delta y^2,$$
$$[y,z]=u(x)+\big(\alpha\delta^{-1} x^2+\gamma_1+\alpha\beta\delta^{-1}\big)y +\big(\alpha\delta^{-1} x+\gamma_2+\alpha\delta^{-1}\big)z,$$
where $\alpha\in\{0,1\}$, $\delta,\beta,\gamma_2,\gamma_1,\gamma_0\in K$ with $\delta\neq 0$, and $u(x)\in K[x]$ has degree at most $4$. 

\item When the characteristic of $K$ is $2$, we get the family with generators $x,y,z$ and relations
$$[x,y]=z, [x,z]=  z + \beta y + ( \gamma x^2 + \gamma x + \zeta ) +\delta y^2,  [y,z]= \gamma z,$$ where $\beta, \gamma, \zeta, \delta\in K$.
\item When the characteristic of $K$ is $2$, we get the family with generators $x,y,z$ and relations
$$[x,y]=z,\quad [x,z]=  \epsilon z +\beta y + (\gamma_3 x^3+ \gamma_2 x^2 + \gamma_1 x + \gamma_0 )$$
and $$[y,z] = u(x)+(\gamma_3 x^2+\gamma_1+\beta\gamma_3)y+(\gamma_3 x+\gamma_2+\epsilon\gamma_3)z+\epsilon y^2,$$
where $u(x)\in K[x]$ has degree at most four and $\beta,\gamma_0,\gamma_1,\gamma_2,\gamma_3,\epsilon\in K$.

\item When the characteristic of $K$ is $2$, we get the family with generators $x,y,z$ and relations
$$[x,y]=z, [x,z]=\epsilon z, [y,z]= u(x)+v(x)(\epsilon y+z)+\epsilon y^2,$$
where $\epsilon\in K, u(x),v(x)\in K[x]$ and with $v$ of degree at most one and $u$ of degree at most four.
\end{enumerate} 
\end{thm}
\begin{remark}
We note that although we are using the word ``classification'' we do not claim that this list is irredundant up to isomorphism, and so the isomorphism problem within and among families is not resolved; moreover, when $K$ is algebraically closed, family (2) can also be realized as an iterated Ore extension by making a change of variables $x'=x, y'\in Ky+Kz, z'=y$ with $y'$ a linear combination of $y$ and $z$ (involving $z$ non-trivially) that makes $y'$ an eigenvector for ${\rm ad}_x$.  In particular, when $K$ is algebraically closed and of characteristic zero, the classification consists precisely of families (1), (3), and (4).
\end{remark}

The outline of this paper is as follows. In \S \ref{sec:prelim} we establish our notation and the cases we study to obtain the classification result, Theorem \ref{thm:main}. In \S\ref{sec:I}--\ref{sec:IIIb} we work through each of the cases and obtain the families given in Theorem \ref{thm:main}. In \S\ref{sec:main} and \S\ref{sec:et} we prove Theorems \ref{thm:main} and \ref{thm:main2} respectively.

\section{Preliminaries for proof of Theorem \ref{thm:main}}\label{sec:prelim}
We let $K$ be a field and let $A$ be a filtered deformation of $K[x,y,z]$.  In particular, $A$ is equipped with an ascending filtration 
$$K=V_0\subseteq V_1\subseteq V_2\subseteq \cdots$$ such that each $V_i$ is finite-dimensional, $A=\bigcup V_n$, and the associated graded algebra $\bigoplus_{n\ge 1} V_n/V_{n-1}$ of $A$ with respect to this filtration is isomorphic to $K[x,y,z]$.  Given a nonzero element $u\in A$, we let ${\rm deg}(u)$ denote the smallest nonnegative integer $n$ such that $u\in V_n$, and call this the degree of $u$.  We take $-\infty$ to be the degree of $0$.  Since the associated graded ring of $A$ is commutative, we have the inequality
$${\rm deg}([u,v])\le {\rm deg}(u)+{\rm deg}(v)-1$$ for $u,v\in A$.  Since the associated graded ring of $A$ is a polynomial algebra in three variables, we have that $A$ has generators $x,y,z$ and $A$ has a $K$-basis given by monomials $x^i y^j z^k$ with $i,j,k\ge 0$.  We let $a,b$, and $c$ denote the respective degrees of $x,y$, and $z$ and we may assume that $a\le b\le c$.  Then since $[x,y]$ has degree at most $a+b-1$, we see that $[x,y]$ is a linear combination of monomials $x^i y^j z^k$ in which either $k=1$ and $i=j=0$; or $j=1$ and $i=k=0$; or $j=k=0$.  In particular, if $z$ occurs with nonzero coefficient in our expression for $[x,y]$ as a linear combination of monomials $x^i y^j z^k$, then we have
$[x,y]=\alpha z+h(x,y)$ for some nonzero $\alpha\in K$ and some polynomial $h=\sum c_{i,j} x^i y^j$, depending only on $x$ and $y$, and so we may replace the generator $z$ by $\alpha z+h(x,y)$ and assume without loss of generality that $[x,y]=z$.
Alternatively, if $z$ does not occur in our expression for $[x,y]$ then there are two possibilities:
\begin{itemize}
\item either $[x,y]=\alpha y+h(x)$ for some nonzero $\alpha\in K$ and some polynomial $h(x)$, in which case, we can replace the generator $y$ with $\alpha y+h(x)$ and the generator $x$ with $\alpha^{-1} x$ and assume that $[x,y]=y$; or
\item $[x,y]=h(x)$ for some polynomial $h(x)$.
\end{itemize}
Thus after making these substitutions we may assume that our algebra has three generators $x,y,z$ with respective degrees $a,b,c$ and $a\le b\le c$ such that one of the following must occur:
\begin{enumerate}
\item[Case I:] $[x,y]=h(x)$ for some $h(x)\in K[x]$;
\item[Case II:] $[x,y]=y$;
\item[Case III:] $[x,y]=z$.
\end{enumerate}
As it turns out, Case III is the most difficult case to study, with the other cases being considerably simpler.  We thus classify filtered deformations in Cases I and II first and divide Case III into two subcases:
\begin{enumerate}
\item[Case III(a):] $[x,y]=z$, and $b>2a-2$.
\item[Case III(b):] $[x,y]=z$ and $b\le 2a-2$.
\end{enumerate}

\section{Filtered deformations in Case I}\label{sec:I}
We recall that we have a $K$-algebra generated by elements $x,y,z$ with respective degrees $a,b,c$ and $a\le b\le c$ and such that $[x,y]=h(x)$.  This means that 
the subalgebra $B$ of $A$ generated by $x$ and $y$ is a skew polynomial ring of derivation-type $K[x][y;\delta]$ in which $\delta(x)=h(x)$.  Now since $x^i y^j z^k$ has degree strictly larger than $a+c-1$ if $k\ge 1$ and $i+j\ge 1$, we see that 
$[x,z]$ can only involve the monomials $z$ and monomials $x^i y^j$, thus $[x,z]=\alpha z+ h(x,y)$ for some $\alpha\in K$ and some $h(x,y)\in B$.  Similarly, $[y,z]$ can only involve monomials of the form $x^i z$ with $i\ge 0$ and monomials of the form $x^i y^j$ and so $[y,z]=u(x) z+ g(x,y)$ for some $u(x)\in K[x]$ and some $g(x,y)\in B$.  
Then it is straightforward to show that $A=K[x][y;\delta][z;\sigma,\mu]$ with $\sigma$ a unipotent automorphism given by $x\mapsto x-\alpha, y\mapsto y-u(x)$ and $\mu$ is a $\sigma$-derivation defined by $\mu(x)=h(x,y), \mu(y)=g(x,y)$. This is an iterated Ore extension of unipotent type and hence part of family (1) in Theorem \ref{thm:main}.
\section{Filtered deformations in Case II}\label{sec:II}
Here we have a subalgebra $B$ of $A$ generated by $x$ and $y$ and the relation $[x,y]=y$ gives that $B=K[y][x;\delta]$ with $\delta(y)=y$. As in Case I, we have that
$[x,z] =\alpha z+ h(x,y)$ for some $\alpha\in K$ and some $h(x,y)\in B$ and we have $[y,z]=u(x)z + g(x,y)$ for some $u(x)\in K[x]$ and some $g(x,y)\in B$.
If we let $\sigma$ be the unipotent automorphism of $B$ given by $\sigma(x)=x-\alpha, \sigma(y)=y-u(x)$ and let $\mu$ be the $\sigma$-derivation defined by $\mu(x)=h(x,y)$ and $\mu(y)=g(x,y)$, then $A=B[z;\sigma,\mu]$ is an iterated Ore extension.  The constraints on $g,h,u,\alpha$ needed for our algebra to be an iterated Ore extension (i.e., we need $\delta(xy-yx-y)=0$) are precisely those that ensure that the Jacobi identity holds and our algebras are covered by (1) in Theorem \ref{thm:main}.  
\section{Filtered deformations in Case III(a)}\label{sec:IIIa}
In this case $c\le a+b-1$, so $[x,z]$ has degree at most $2a+b-2$. In particular, since this is less than $2b$ we see that we cannot have any monomials $x^i y^j z^k$ with $j\ge 2$ or with $k\ge 1$ unless $(i,j,k)=(0,0,1)$. Thus 
$[x,z]= u(x) +v(x) y + \alpha z$ for some polynomials $u(x), v(x)\in K[x]$ and $\alpha\in K$. Moreover, $v(x)$ is at most linear. Similarly, 
$[y,z]$ has degree at most $b+c-1\le 2b+a-2$ and so $[y,z]$ can only involve monomials $y^2$ and monomials of the form $x^i, x^j y, x^k z$.
Thus we see in this case there are constants $\alpha,\beta,\gamma,\delta \in K$ and polynomials $u(x), f(x), g(x), h(x)$
such that we have relations of the form
$$[x,y]=z, \quad [x,z] = u(x) + (\beta x+\gamma)y + \alpha z, \quad [y,z]=f(x)+g(x)y + h(x)z+\delta y^2.$$

\subsection{Subcase 1: $(\beta,\gamma)\neq (0,0)$.}
In this case, we can write $u(x) = (\beta x + \gamma) v(x) + \zeta$ for some polynomial $v(x)$ and some $\zeta\in K$, where $\zeta=0$ if $\beta=0$. We can make a substitution 
$$x' = x, \quad y' = y+v(x), \quad z' = z,$$ 
and assume that
$$[x,y]=z, \quad [x,z] = \zeta + (\beta x+\gamma) y + \alpha z, \quad [y,z] = f(x) + g(x) y+h(x)z+\delta y^2.$$
We note that the general form of $[y,z]$ remains unchanged under this substitution, because we can always choose our degree so that the degree of $v(x)$ is smaller than ${\rm deg}(y)=b$ and we can ensure that our substitution algebra preserves the property of being a filtered deformation of a polynomial algebra.
Then the Jacobi identity constraint gives
$$[[x,y],z] = [[x,z],y]+[x,[y,z]].$$
Since $[[x,y],z] = [z,z] = 0$, using our relations we see we must have
$$\beta zy + \alpha [z,y] + g(x) z + h(x) [x,z] + \delta yz + \delta zy = 0,$$
which simplifies to
\begin{equation}
\label{eq:bdg}
(\beta+2\delta) yz - (\alpha +\beta+\delta) [y,z] + g(x) z + h(x)(\zeta + (\beta x + \gamma)y + \alpha z) = 0.
\end{equation}
Then considering the coefficient of $yz$ on both sides, we see that $\beta=-2\delta$. 
Considering the coefficient of $y^2$ on both sides, we see that $(\alpha+\beta+\delta) \delta=(\alpha-\delta)\cdot \delta=0$, so either $\delta=0$ or $\alpha=\delta$.  
In particular, if $\delta=0$, then $\beta=0$ and so our commutators simplify to:
$$[x,y]=z, \quad [x,z] = \gamma y +\alpha z, \quad [y,z]=f(x)+g(x)y + h(x) z.$$
The Jacobi identity $0 = [[x,z],y] + [x,[y,z]]$ now yields:
$$0 = [\gamma y + \alpha z, y] + [x, f(x) + g(x)y + h(x)z],$$
which expands to
$$0 = \alpha [z,y] + g(x)[x,y] + h(x)[x,z].$$
Substituting the commutators gives:
$$0 = -\alpha [y,z] + g(x)z + h(x)(\gamma y + \alpha z),$$
or equivalently, moving the $\alpha[y,z]$ term to the left hand side:
$$\alpha (f(x) + g(x)y + h(x)z) = \gamma h(x)y + (g(x) + \alpha h(x))z.$$
Comparing coefficients as modules over $K[x]$ then gives
$$\alpha h(x) = g(x) + \alpha h(x) \implies g(x) = 0.$$
Looking at the $y$ coefficient gives $\alpha g(x) = \gamma h(x)$. Since $g(x) = 0$, this becomes $0 = \gamma h(x)$. By the assumption of this subcase, $(\beta,\gamma)\neq (0,0)$. Having already shown $\beta=0$, it follows that $\gamma \neq 0$, which forces $h(x)=0$. 

Comparing the $x$-polynomial coefficients gives $\alpha f(x) = 0$. Thus, our commutators simplify to:
$$[x,y] = z, \quad [x,z] = \gamma y + \alpha z, \quad [y,z] = f(x),$$
with the constraint that either $\alpha = 0$ or $f(x) = 0$.
If $f(x)=0$, then our algebra is an iterated Ore extension $K[y,z][x;\delta]$ where $\delta(y)=z, \delta(z)=\gamma y$ and otherwise we get an algebra
\begin{equation}
[x,y]=z, \qquad [x,z]=\gamma y, \qquad [y,z]=f(x),
\end{equation}
which is family (2) in the statement of Theorem \ref{thm:main}.  Thus when $\delta=0$ we see this case yields families from (1) and (2) in Theorem \ref{thm:main}.

When $\delta\neq 0$, we have $\alpha=\delta, \beta=-2\delta$ so our commutators become
$$[x,y]=z, \quad [x,z] = \zeta + (-2\delta x+\gamma) y + \delta z, \quad [y,z] = f(x) + g(x) y+h(x)z+\delta y^2.$$
If the characteristic of $K$ is not $2$ then $2\delta\neq 0$.  We can in this case make a change of variables
$x'=\delta x-\gamma/2$, $y' = y$, $z'=\delta z$ and our relations are now of the form
$$[x,y]=z, \quad [x,z] = -2xy + z+\zeta, \quad [y,z] = f(x) + g(x) y+h(x)z+y^2.$$
Since $\alpha+\beta+\delta=0$ and since $\beta x+ \gamma$ is nonzero, looking at the $x$-polynomial coefficients of $y$ on both sides of Equation (\ref{eq:bdg}) gives
$h(x)=0$.  Then, looking at the $x$-polynomial coefficients of $z$ gives $g(x)+\alpha h(x)=0$ and since $\alpha=\delta\neq 0$, we see $g(x)=0$, so we get 
the commutators
$$[x,y]=z,\quad [x,z] =z-2xy +\zeta, \quad [y,z]=y^2+f(x),$$ which is family (4) from Theorem \ref{thm:main}.

Finally, if the characteristic of $K$ is two and $\delta\neq 0$ then $\beta=-2\delta=0$, so by assumption $\gamma\neq 0$ and since $\beta=0$, by the remarks above $\zeta=0$ and so our commutators are given by 
$$[x,y]=z, [x,z] = \gamma y + \delta z, [y,z] = \delta y^2 + f(x) + g(x) y + h(x) z.$$
Then the Jacobi identity gives the constraints
$$0=[[x,y],z]=[x,[y,z]]+[[x,z],y] = (\delta [y,z] + g(x) z + h(x) [x,z])+\delta [z,y])$$ and so
$$g(x)z + h(x) (\gamma y + \delta z)=0.$$ Since $\gamma\neq 0$ by assumption, we see $h(x)=g(x)=0$ and so we see that the algebras we obtain are from family (8) with $\gamma_3=\gamma_2=\gamma_1=\gamma_0 =0$ in this case.


\subsection{Subcase 2: $\beta=\gamma=0$}
In this case, our relations are
$$[x,y]=z, \quad [x,z] =  u(x) +  \alpha z, \quad [y,z] = f(x) + g(x) y+h(x)z.$$
The Jacobi identity gives the constraint
$[x,[y,z]]=[[x,z],y]$.  The coefficients of terms of the form $x^iy$ on the left side can be checked to be $0$, whereas the contribution from these terms on the right side is
$\alpha g(x)y$, which can be observed by noting that 
the Hasse-Schmidt expansion for the bracket $[u(x), y] = \sum_{n=1}^{\infty} (-1)^{n-1} u^{[n]}(x) \text{ad}_x^n(y)$ shows that $[y,u(x)]$ involves no monomials of the form $x^iy$ when we put our expression in normal form and $[\alpha z,y]$ contributes $\alpha g(x)y$.  It follows that $\alpha g(x)=0$.

If $g(x)=0$ then our algebra is given by the relations 
$$[x,y]=z, \quad [x,z] =  u(x) +  \alpha z, \quad [y,z] = f(x) +h(x)z$$ and so it is an iterated Ore extension of unipotent type of the form $K[x][z;\sigma_1,\delta_1][y;\delta_2]$.

On the other hand, if $\alpha=0$ then our algebra is of the form
$$[x,y]=z, \quad [x,z] =  u(x), \quad [y,z] = f(x) + g(x) y+h(x)z,$$ which is again an iterated Ore extension of unipotent type of the form
$K[x][z;\delta_1][y;\sigma_2,\delta_2]$.


\section{Filtered deformations in Case III(b)} \label{sec:IIIb}
Here since $b\le 2a-2$ and $z=[x,y]$ we see that $c\le a+b-1\le 3a-3$.  Thus
$[x,z]$ has degree at most $2a+b-2\le 4a-4$ and $[y,z]$ has degree at most $a+2b-2\le 5a-6$.  In particular, $[x,z]$ is a linear combination of monomials $x^i y^j z^k$ from
$$\{1,x,x^2,x^3, y, xy, y^2, z\}$$ and 
$[y,z]$ is a linear combination of monomials $x^i y^j z^k$ from
$$\{1,x,x^2,x^3, x^4, y, xy, x^2y, y^2, z, xz\}.$$  
So we can write
\begin{equation}
[x,z]= \alpha z + (\beta_1 x +\beta_0) y + (\gamma_3 x^3+ \gamma_2 x^2 + \gamma_1 x + \gamma_0 ) +\delta y^2
\end{equation}
and
\begin{equation}
[y,z]= u_0(x) + u_1(x) y + u_2(x) z + \epsilon y^2 
\end{equation}
where $u_0, u_1, u_2$ have degrees at most $4, 2$, and $1$ respectively.
We write $$u_0(x)=\textstyle\sum_{i=0}^4 c_ix^i,\ u_1(x)=e_2x^2+e_1x+e_0,\ u_2(x)=f_1x+f_0.$$
This is the most involved case and we use six subcases to deal with this case, based on whether the values $\alpha, \beta_1$, $\delta,$ and $\epsilon$ are zero or not.
For the analysis that follows, we will need these formulas:

\begin{equation} \label{eq:xxy}
[x^2,y] = xz+zx = 2xz - [x,z],
\end{equation}
\begin{equation} \label{eq:xxxy}
 [x^3, y] = x^2 z + xzx + zx^2 = 3x^2z - 3x[x,z] + [x,[x,z]].
 \end{equation}

\subsection{Subcase 1: $\delta=\epsilon=0$.}
In this case, the commutator brackets take the form of those given in Case III(a), and so we obtain no new deformations in this case.


\subsection{Subcase 2: $\delta=0$, $\epsilon \neq 0$.}
In this case the constraints given by the Jacobi identity give in particular $\alpha=\epsilon, \beta_1=-2\epsilon$.  Thus if the characteristic of $K$ is not $2$, $\beta_1=-2\epsilon \neq 0$.
By the division algorithm $$\gamma_3 x^3+ \gamma_2 x^2 + \gamma_1 x + \gamma_0= (-2\epsilon x+\beta_0) q(x) + \zeta$$ for some $q(x)$ of degree at most $2$ and some $\zeta\in K$.  By making a substitution $x'=x-\beta_0\epsilon^{-1}/2, y'=y-q(x), z'=z$, we see that the general form of the $[y,z]$ commutator remains unchanged so we get
$$[x,y]=z, [x,z]=\epsilon z -2\epsilon xy +\zeta, [y,z]=\epsilon y^2 +u_0(x)+u_1(x)y+u_2(x)z.$$
The remaining constraints from the Jacobi identity give:
$$ e_2=0, e_1=-\epsilon f_1, e_0=-\epsilon f_0, 2 f_1=0, -2\epsilon f_0=0, \zeta f_0 =0$$

Thus if the characteristic of $K$ is not $2$, this immediately gives us the family
$$[x,y]=z, [x,z]=\epsilon z -2\epsilon xy +\zeta, [y,z]=\epsilon y^2 +u_0(x),$$ which appears as family (4) in Theorem \ref{thm:main}.

If the characteristic of $K$ is $2$, then we have $\alpha=\epsilon$ and $\beta_1=2\epsilon=0$ and so our commutators take the form
$$[x,z]= \epsilon z +\beta_0 y + (\gamma_3 x^3+ \gamma_2 x^2 + \gamma_1 x + \gamma_0 )$$
and
$$[y,z]= u_0(x) + u_1(x) y + u_2(x) z + \epsilon y^2.$$
We find $$[[x,z],y]=\epsilon [y,z] + (\gamma_3 x^2+\gamma_1)z + (\gamma_3 x+\gamma_2)[x,z]+\gamma_3[x,[x,z]]$$
while $$[[y,z],x]= u_1(x) z + u_2(x) [x,z]+ \epsilon [y,z].$$
Since $[x,[x,z]]= \epsilon [x,z] + \beta_0 z$ and since $z$ and $[x,z]$ span a $2$-dimensional left $K(x)$-vector space when either $\beta_0\neq 0$ or $(\gamma_3,\gamma_2,\gamma_1,\gamma_0)\neq (0,0,0,0)$, when one of these conditions holds the Jacobi identity gives
$$ (\gamma_3 x^2+\gamma_1+\beta_0\gamma_3)=u_1(x)$$ and
$$(\gamma_3 x+\gamma_2+\epsilon\gamma_3) = u_2(x).$$
Thus we get the commutators 
$$[x,y]=z,\quad [x,z]=  \epsilon z +\beta_0 y + (\gamma_3 x^3+ \gamma_2 x^2 + \gamma_1 x + \gamma_0 )$$
and $$[y,z] = u_0(x)+(\gamma_3 x^2+\gamma_1+\beta_0\gamma_3)y+(\gamma_3 x+\gamma_2+\epsilon\gamma_3)z.$$
This is family (8) in Theorem \ref{thm:main}.

On the other hand, if $\beta_0=\gamma_3=\gamma_2=\gamma_1=\gamma_0=0$ we have the commutators
$$[x,y]=z, [x,z] =\epsilon z, [y,z] = u_0(x) + u_1(x) y + u_2(x) z + \epsilon y^2.$$
and the Jacobi identity gives
$$\epsilon [y,z] = [y,[x,z]] = [x,[y,z]] = u_1(x) z + u_2(x) \epsilon z + \epsilon [z,y],$$ and so we require that $u_1(x)=\epsilon u_2(x)$ and so we have the family
$$[x,y]=z, [x,z]=\epsilon z, [y,z]= u_0(x)+u_2(x)(\epsilon y+z)+\epsilon y^2.$$
This is family (9) in Theorem \ref{thm:main}.
 
\subsection{Subcase 3: $\alpha\beta_1\delta\neq 0$}
In this case, by first making a substitution of the form $x'=\alpha^{-1} x, y'=\alpha y, z'=z$, we can assume $\alpha=1$. Next by making a substitution of the form $x'=x+\beta_1^{-1}\beta_0$, $y'=y, z'=z$, we may assume that $\beta_0=0$.  Finally, we can make a substitution $x'=x, y'=y-\epsilon \delta^{-1} x$ to assume that $\epsilon =0$.  Thus we may assume that our commutators are given by
$$[x,y]=z,\quad [x,z]= z + \beta_1 x y + (\gamma_3 x^3+ \gamma_2 x^2 + \gamma_1 x + \gamma_0 ) +\delta y^2$$
and
$$[y,z]= u_0(x) + u_1(x) y + u_2(x) z.$$
Then the Jacobi identity constraint from $0=[[x,y],z]$ gives $[[x,z],y]= -[x,[y,z]]$. 
We compute $[[x,z],y]$ using Equations (\ref{eq:xxy}) and (\ref{eq:xxxy}) and find
$$[[x,z],y] = [z,y] + \beta_1 zy + \gamma_3 (3x^2 z-3x[x,z]+[x,[x,z]]) + \gamma_2 (2xz -[x,z]) + \gamma_1 z.$$
While $$-[x,[y,z]]= -u_1(x) z - u_2(x) (z+\beta_1 xy +(\gamma_3 x^3+ \gamma_2 x^2 + \gamma_1 x + \gamma_0 ) +\delta y^2).$$
Simplifying and comparing coefficients we obtain the following constraints:
$$\beta_1 =- 2\delta\gamma_3, \quad f_1 = 3\gamma_3, \quad  f_0 = \gamma_2 - \gamma_3 $$
 $$e_2 = -3\gamma_3, \quad e_1 = 3\gamma_3 M - 2\gamma_2 - \gamma_3\beta_1, 
    e_0 = (\gamma_2-\gamma_3)M - \gamma_1,$$
    and $$M e_2  = M e_1 = M e_0=0,\quad M \cdot u_0(x) =0,$$
   where $ M = 1-\delta\gamma_3$.
   
In particular, since $\beta_1\neq 0$ is a multiple of $2$, the characteristic of $K$ cannot be equal to $2$ in this subcase, so we assume for the remainder of this subcase that $K$ does not have characteristic $2$.

If $M\neq 0$ and the characteristic of $K$ is not $3$, then the constraints give $u_0(x)=e_2=e_1=e_0=\beta_1=\gamma_3=\gamma_2=\gamma_1=f_1=f_0=0$, and this is not allowed since we are assuming that $\beta_1\neq 0$.

If $M\neq 0$ and the characteristic of $K$ is $3$, we get constraints:
$$u_0(x)=e_2=e_1=e_0=f_1=0, \beta_1=\delta \gamma_3, f_0=\gamma_2-\gamma_3, \gamma_2=\delta \gamma_3^2, \gamma_1=(\gamma_2-\gamma_3)M.$$
This gives us the family
\begin{equation}
[x,y]=z,\quad [x,z]= z +\delta\gamma_3xy + \delta y^2 + (\gamma_3 x^3 +\delta \gamma_3^2 x^2-\gamma_3 (1-\delta \gamma_3)^2 x+\gamma_0),
\end{equation}
and
\begin{equation} [y,z]= (\gamma_3^2\delta -\gamma_3)z,
\end{equation}
which is the family from item (5) in Theorem \ref{thm:main}.



If $M=0$ then we get $\gamma_3=\delta^{-1}$, $\beta_1=-2$, $f_1=3\delta^{-1}$, $f_0=\gamma_2-\delta^{-1}$, $e_2=-3\gamma_3$, $e_1=-2\gamma_2-\delta^{-1}\beta_1$, $e_0=-\gamma_1$
%$e_2=-3\gamma_3$, $e_1=-2\gamma_2+2\delta^{-1}$, $e_0=-\gamma_1$

and our
commutator brackets simplify to 
\begin{equation}
[x,y]=z,\quad [x,z]= z -2xy + \delta y^2 + (\delta^{-1} x^3 +\gamma_2 x^2+\gamma_1 x+\gamma_0),
\end{equation}
and
\begin{equation} [y,z]= u_0(x) - (\gamma_1+(2\gamma_2 -2\delta^{-1})x + 3\delta^{-1}x^2) y +
(\gamma_2-\delta^{-1} +3\delta^{-1} x) z,
\end{equation}
which is family (3) from Theorem \ref{thm:main} with $\alpha=1$.




\subsection{Subcase 4: $\delta\neq 0$, $\alpha=0$, $\beta_1
\neq 0$}
As in Subcase 3, after the substitutions $x'=x+\beta_1^{-1}\beta_0$ and
$y'=y-\epsilon\delta^{-1}x$ we may assume $\beta_0=\epsilon=0$; we cannot, however,
normalize $\alpha$, which is now zero. Thus we may assume that our commutators are given by
$$[x,y]=z,\quad [x,z]= \beta_1 x y + (\gamma_3 x^3+ \gamma_2 x^2 + \gamma_1 x + \gamma_0 ) +\delta y^2$$
and
$$[y,z]= u_0(x) + u_1(x) y + u_2(x) z.$$
The Jacobi identity gives $[[x,z],y]=-[x,[y,z]]$, and using Equations
(\ref{eq:xxy}) and (\ref{eq:xxxy}) we find
$$[[x,z],y] = \beta_1 zy + \gamma_3 (3x^2 z-3x[x,z]+[x,[x,z]]) + \gamma_2 (2xz -[x,z]) + \gamma_1 z,$$
while
$$-[x,[y,z]]= -u_1(x) z - u_2(x) \big(\beta_1 xy +(\gamma_3 x^3+ \gamma_2 x^2 + \gamma_1 x + \gamma_0 ) +\delta y^2\big).$$
Comparing coefficients we obtain
$$\beta_1 =- 2\delta\gamma_3, \quad f_1 = 3\gamma_3, \quad  f_0 = \gamma_2,$$
$$e_2 = -3\gamma_3, \quad e_1 = 3\gamma_3 N - 2\gamma_2 - \gamma_3\beta_1, \quad e_0 = \gamma_2 N - \gamma_1,$$
$$N e_2  = N e_1 = N e_0 = N\cdot u_0(x) = 0, \qquad N=-\delta\gamma_3.$$
Since $\beta_1=-2\delta\gamma_3
\neq 0$ we have $\gamma_3
\neq 0$ and $\mathrm{char}(K)
\neq 2$.

In particular, $N=-\delta\gamma_3
\neq 0$, so the last line forces
$$e_2=e_1=e_0=0,\qquad u_0(x)=0.$$
Combining $e_2=0$ with $e_2=-3\gamma_3$ gives $3\gamma_3=0$, and since $\gamma_3
\neq 0$
this is possible only when $\mathrm{char}(K)=3$. Hence there are no deformations in this
subcase unless $K$ has characteristic $3$.

When the characteristic of $K$ is three, the constraints simplify: $e_1=0$ gives
$\gamma_2=\delta\gamma_3^2$; $e_0=0$ gives $\gamma_1=\gamma_2 N=-\delta^2\gamma_3^3$; and
$f_1=3\gamma_3=0$, $f_0=\gamma_2=\delta\gamma_3^2$. Thus $u_0=u_1=0$ and
$u_2=\delta\gamma_3^2$, giving the family
$$[x,y]=z,\quad [x,z]= \delta\gamma_3\, xy + \delta y^2 + \big(\gamma_3 x^3 +\delta \gamma_3^2 x^2-\delta^2\gamma_3^3\, x+\gamma_0\big),$$
$$[y,z]= \delta\gamma_3^2\, z,$$
with $\delta
\neq 0$, $\gamma_3
\neq 0$, $\gamma_0\in K$. This is family (5) of
Theorem~\ref{thm:main} with $\alpha=0$.




%\subsection{Subcase 4: $\delta\neq 0$, $\alpha=0, \beta_1\neq 0$}
%As in the preceding Subcase, after making substitutions we may assume that $\beta_0=\epsilon=0$.  Using the Jacobi identity, we find that $\beta_1=2\alpha=0$ and so there are no deformations in this subcase.



\subsection{Subcase 5: $\delta\neq 0$, $\alpha\neq 0, \beta_1=0$}


As in Subcase 3, after the substitutions $x'=\alpha^{-1}x,\ y'=\alpha y$ and
$y'=y-\epsilon\delta^{-1}x$ we may assume $\alpha=1$ and $\epsilon=0$, so that
$$[x,y]=z,\quad [x,z]= z + \beta_0 y + (\gamma_3 x^3+ \gamma_2 x^2 + \gamma_1 x + \gamma_0 ) +\delta y^2,\quad [y,z]= u_0(x) + u_1(x) y + u_2(x) z.$$
The Jacobi identity gives $[[x,z],y]=-[x,[y,z]]$, and using Equations
(\ref{eq:xxy}) and (\ref{eq:xxxy}) the coefficient of $yz$ yields
$$2\delta\gamma_3=0.$$

Suppose first that the characteristic of $K$ is not two.
Then $\gamma_3=0$, and comparing the
remaining coefficients gives $\gamma_2=\gamma_1=0$ together with
$u_0(x)=u_1(x)=u_2(x)=0$. Hence $[y,z]=0$ and
$$[x,y]=z,\quad [x,z]=z+\beta_0 y+\gamma_0+\delta y^2,\quad [y,z]=0,$$
so that $A=K[y,z][x;\sigma,\delta]$ is an iterated Ore extension of unipotent
type and lies in family (1). We therefore assume for the
remainder of this subcase that the characteristic of $K$ is two.

Under this assumption the constraint $2\delta\gamma_3=0$ is vacuous, and solving the
Jacobi identity $[[x,z],y]=[[y,z],x]$ we obtain the conditions
$$f_0=\gamma_2+\gamma_3,\ f_1=\gamma_3,\ e_2=\gamma_3,\ e_1=\gamma_3+\delta\gamma_3^2,\ e_0=\gamma_1+\gamma_2+\gamma_3+\beta_0\gamma_3+\delta\gamma_3(\gamma_2+\gamma_3)$$
along with $(\delta\gamma_3+1)Z=0$ for all $Z\in\{\gamma_3,\gamma_1+\gamma_2,c_0,c_1,c_2,c_3,c_4\}$.
Thus when $\gamma_3=\delta^{-1}$ our commutators simplify to
$$[x,y]=z,\quad [x,z]=z + \beta_0 y + (\delta^{-1} x^3+ \gamma_2 x^2 + \gamma_1 x + \gamma_0 ) +\delta y^2,$$
and
$$[y,z] = u_0(x) + (\delta^{-1} x^2+\gamma_1+\beta_0\delta^{-1})y + (\delta^{-1} x+\gamma_2+\delta^{-1})z,$$
which is part of family (6) in Theorem \ref{thm:main}.
When $\gamma_3
\neq \delta^{-1}$ we get the additional constraints
$\gamma_3=\gamma_1+\gamma_2=u_0(x)=0$, and so our commutators become
$$[x,y]=z,\quad [x,z]=  z + \beta_0 y + ( \gamma_1 x^2 + \gamma_1 x + \gamma_0 ) +\delta y^2,\quad [y,z]= \gamma_1 z,$$
which is family (7) in Theorem \ref{thm:main}.




\subsection{Subcase 6: $\delta\neq 0$, $\alpha=\beta_1=0$}

In this case, as in Subcase 2, we may make a substitution and assume that $\epsilon=0$.
Thus our commutator brackets can be assumed to be of the form
 $$[x,y]=z,\quad [x,z]=   \beta_0 y + (\gamma_3 x^3+ \gamma_2 x^2 + \gamma_1 x + \gamma_0 ) +\delta y^2$$
and
$$[y,z]= u_0(x) + u_1(x) y + u_2(x) z.$$

We compute the coefficient constraints given by the Jacobi identity and obtain the following constraints:
$$f_0=\gamma_2, f_1=3\gamma_3, e_2=-3\gamma_3, e_1=3\delta\gamma_3^2-2\gamma_2, e_0=\delta\gamma_2\gamma_3-\beta_0\gamma_3-\gamma_1,$$
$$2\gamma_3 = 3\gamma_3^2=0 \implies \gamma_3=0$$
Thus our commutators simplify to
\begin{equation}
[x,y]=z,\quad [x,z]=\beta_0 y + \gamma_2 x^2 + \gamma_1 x + \gamma_0 +\delta y^2, \quad [y,z]=u_0(x) + (-2\gamma_2 x-\gamma_1)y +  \gamma_2 z.
\end{equation}
In the case that the characteristic of $K$ is equal to $2$, this simplifies to 
$$[x,y]=z, \quad [x,z]= \beta_0 y +  \gamma_2 x^2 + \gamma_1 x + \gamma_0 +\delta y^2,$$ $[y,z]=u_0(x)+\gamma_1y +\gamma_2 z$, which  
is just family (6) in Theorem \ref{thm:main}.

If the characteristic of $K$ is not $2$ then we can make a substitution $x'=x, y'=y+\delta^{-1}\beta_0/2, z'=z$ and our family simplifies to
\begin{equation}
[x,y]=z,\quad [x,z]= \gamma_2 x^2 + \gamma_1 x + \gamma_0 +\delta y^2, \quad [y,z]=u_0(x) + (-2\gamma_2 x-\gamma_1)y +  \gamma_2 z.
\end{equation}
This is a subfamily of the family (3) in Theorem \ref{thm:main}.

\section{Proof of Theorem \ref{thm:main}}\label{sec:main}
We can now quickly prove Theorem \ref{thm:main}. We note that it is straightforward to show computationally that each of the families (1)--(9) in Theorem \ref{thm:main} are filtered deformations of a polynomial ring in three variables.  To do this, we use the relations to put each of $[[x,y],z], [[x,z],y]$, and $[x,[y,z]]$ in normal form as a linear combination of monomials $x^i y^j z^k$. One then checks that the Jacobi identity holds via computer calculation. One next uses the filtered property to put a degree on $x$, $y$, $z$ that makes 
${\rm deg}([u,v]) < {\rm deg}(u)+{\rm deg}(v)$ for $u,v\in \{x,y,z\}$.  Next, using a degree-lexicographic ordering on words in $x,y,z$ with $x>y>z$, Bergman's diamond lemma (see specifically his proof of the PBW theorem) \cite{Berg} shows $xy, xz, yz$ are leading monomials of a Gr\"obner basis for the ideal generated by relations and so our algebra is indeed a filtered deformation of a polynomial ring. On the other hand, the analysis performed in studying the algebras that arise in Cases I, II, III(a), and III(b) shows that every filtered deformation of a polynomial ring is necessarily of one of these forms.  

\section{Etingof's question for three-variable polynomial rings}\label{sec:et}
We can now show that in positive characteristic a filtered deformation of a polynomial ring in three variables is PI.  To do this, we use the classification provided in Theorem \ref{thm:main}. We need a criterion that ensures certain filtered deformations are PI.  To obtain this we first prove a basic lemma.


\begin{lemma} Let $F$ be an algebraically closed field, let $A$ be an $F$-algebra having an ascending filtration by finite-dimensional $F$-vector spaces whose associated graded ring is a finitely generated graded noetherian domain, and let $D$ be the quotient division algebra of $A$.  Then all subfields $K$ of $D$ containing $F$ are finitely generated as extensions of $F$.
\label{lem:vamos}
\end{lemma}
\begin{proof}
Let $B$ be the associated graded of $A$ and let $K$ be a subfield of $D$ containing $F$.  Then by an argument of de Jong appearing in \cite[Theorem 5.1]{ASZ}, $B\otimes_F K$ is noetherian.  In particular, since $B\otimes_F K$ is an associated graded ring of $A\otimes_F K$, $A\otimes_F K$ is noetherian.  Since $D\otimes_F K$ is a localization of $A\otimes_F K$, $D\otimes_F K$ is noetherian. Since $D$ is free over $K$, $D\otimes_F K$ is a free left and right $K\otimes_F K$-module and so a faithful flatness argument gives that $K\otimes_F K$ is noetherian.  It follows from \cite{Vam} that $K$ is finitely generated over $F$.
\end{proof}
The following result is essentially the argument from \cite{Bell}, which shows that filtered deformations of integral domains of Krull dimension two are PI.  This says that the same result holds for filtered deformations of polynomial rings in three variables when they have a non-trivial centre.  Intuitively, this makes sense, as one can invert the nonzero central elements and work over a larger base field and one then obtains a new domain of Gelfand-Kirillov dimension two that is a filtered deformation of an algebra of Krull dimension two.  The only subtleties in invoking the proof from \cite{Bell} are that it requires an algebraically closed base field and one cannot guarantee that the associated graded ring is an integral domain after extension of scalars (or even before extension for that matter!). Nevertheless, it is reasonably straightforward to get around these issues in the setting we are interested in.  

\begin{prop}
\label{prop:invoke}
Let $F$ be an algebraically closed field of positive characteristic $p>0$ and let $A$ be a filtered deformation of a three-variable polynomial ring over $F$.  If $A$ has a central element $\Theta \in A\setminus F$, then $A$ is PI.
\end{prop}
\begin{proof} This proof follows the strategy from \cite{Bell}.  We let $x,y,z$ be generators for $A$ such that $\bar{x},\bar{y},\bar{z}$ generate the associated graded ring of $A$.  Given $u\in A$, we let $\bar{u}$ denote its image in the associated graded ring. 

Then the image of $\Theta$ in the associated graded ring of $A$ is some nonzero homogeneous polynomial $\bar{\Theta}:=h(\bar{x},\bar{y},\bar{z})$.  
We let $a,b,c$ denote the degrees of $x,y$, and $z$ respectively.  Then we can put a degree-lexicographic order on monomials in $\bar{x},\bar{y},\bar{z}$ by ordering first by total degree and then lexicographically using an order $\bar{x}>\bar{y}>\bar{z}$. 

We let $\bar{x}^{\alpha} \bar{y}^{\beta} \bar{z}^{\gamma}$ denote the degree lexicographically greatest monomial that occurs in $\bar{\Theta}$ with nonzero coefficient. After possibly permuting $x$, $y$, and $z$ we may assume that $\alpha>0$ and that this monomial is still degree lexicographically greatest.  

Then by considering images in the associated graded ring and using the fact that $\bar{x}^{\alpha} \bar{y}^{\beta} \bar{z}^{\gamma}$ is the leading monomial of a Gr\"obner basis for the ideal $(\bar{\Theta})$ in $F[\bar{x},\bar{y},\bar{z}]$, we see that $\Theta^i y^j z^k x^s$ with $i,j,k\ge 0$ and $s<\alpha$ spans $A$ as a left $F[\Theta,y]$-module.  We now let $K$ denote the subfield of ${\rm Frac}(A)$ generated by $\Theta$ and $y$.  Then since $\bar{\Theta}=h$ has monomials involving $x$, $\bar{\Theta}$ and $\bar{y}$ are algebraically independent over $F$, and so the field $K$ has transcendence degree $2$ and $KV_n$ is spanned by monomials $z^k x^s$ with $s<\alpha$ such that $k{\rm deg}(z)+s {\rm deg}(x)\le n$.  In particular, 
${\rm dim}_K KV_n/KV_{n-1}=O(1)$ since our basis for $KV_n$ can have at most $\alpha$ elements of degree exactly $n$.  By telescoping, for fixed $r\ge 1$ we have
${\rm dim}_K KV_n/KV_{n-r}=O(1)$ as $n\to\infty$.  

On the other hand, if we let $\Omega_n$ denote the $K$-span of elements $u$ in $A$ such that $[y,u]\in V_n$, then we claim that 
\begin{equation}
\limsup_{n\to \infty} {\rm dim}_K(\Omega_n) - {\rm dim}_K(KV_n) =\infty.\label{eq:inf}\end{equation}
The reason for this is that if we let $L$ be such that ${\rm deg}(y)<p^L$, then if we let $N>L$ and consider the natural numbers
$$a_N:=p^N, a_{N-1}= p^N-p^{N-1}, a_{N-2} = p^N-p^{N-1}-p^{N-2}, \ldots , a_L=p^N-p^{N-1}-\cdots - p^L,$$ then
for $i\in \{L,\ldots ,N\}$, we have $a_i= a_i' p^i$ for some integer $a_i'$ and so
$$[y,z^{a_i}] = -{\rm ad}_{z^{a_i}}(y) = -{\rm ad}_{z^{a_i'}}^{p^i}(y).$$ Since each application of ${\rm ad}_{z^{a_i'}}$ raises the degree by at most $a_i'{\rm deg}(z)-1$, we see that
 $[y,z^{a_i}]$ has degree at most $${\rm deg}(z) a_i - p^i + {\rm deg}(y) \le {\rm deg}(z)(p^N- p^{N-1}-\cdots - 2p^i) \le {\rm deg}(z)\left(p^N- \sum_{j=L-1}^{N-1} p^j\right),$$
and  so $$[y,z^{a_i}]\in \Omega_{b(N)},$$ where $b(N):={\rm deg}(z)(p^N-p^{N-1}-\cdots -p^L - p^{L-1})$.

But now this establishes (\ref{eq:inf}), since the sum $Kz^{a_L}+\cdots +Kz^{a_N}+ KV_{b(N)-{\rm deg}(y)+1}$ is direct and by construction contained in $\Omega_{b(N)}$, so
$\Omega_{b(N)}$ has dimension at least $$N-L+1+{\rm dim}_K(KV_{b(N)-{\rm deg}(y)+1}).$$  Since $N$ is arbitrary and since ${\rm dim}_K(KV_{b(N)}/KV_{b(N)-{\rm deg}(y)+1})=O(1)$ we see (\ref{eq:inf}) holds.  

Now the proof of \cite[Corollary 1.6]{Bell} shows that the centralizer, $E$, of $y$ in ${\rm Frac}(A)$ is a division ring that is infinite dimensional over $K$.  By construction $K\subseteq Z(E)$.
Let $K'$ be a maximal subfield of $E$. If $[K':K]<\infty$, then $[K':Z(E)]<\infty$ and so by \cite[Theorem 15.4]{Lam}, $[E:Z(E)]<\infty$. Thus $[E:K]=[E:Z(E)]\cdot [Z(E):K]<\infty$, which is false.  Hence
$[K':K]=\infty$. By Lemma \ref{lem:vamos}, $K'$ is finitely generated as an extension of $F$ and hence as an extension of $K$ and since $[K':K]=\infty$, $K'$ must have transcendence degree at least one over $K$. Hence the transcendence degree of $K'/F$ is strictly greater than $2={\rm trdeg}_F(K)$ and so $K'$ has transcendence degree $\ge 3$ over $F$. Now \cite[Theorem 1.1]{Belldiv} gives that $A$ is PI.
\end{proof}
We are now ready to prove Theorem \ref{thm:main2}. To do this, we show that the algebras we consider have a non-trivial centre and hence we can invoke Proposition \ref{prop:invoke}. In some families in Theorem \ref{thm:main} (particularly families (3) and (4)), we find centres in part via computer search using Sage Math (code can be found in the appendix of the Master's thesis of the second-named author \cite{Li}).  
We prove the existence of non-trivial central elements for these algebras in the following lemmas.  We first need a lemma that allows us to find a non-trivial centre in characteristic zero to deduce the existence of a non-trivial centre in positive characteristic.
\begin{lemma}\label{lem:lift}
Let $R$ be a finitely generated commutative $\mathbb{Z}$-algebra which is an integral domain,
let $F=\operatorname{Frac}(R)$ be its field of fractions, and let $\theta: R\to K$ be a ring homomorphism from $R$ to a field $K$, so that we view $K$ as an $R$-module.  Let $A$ be an associative
$R$-algebra equipped with an ascending filtration
$$R=A_0\subseteq A_1\subseteq A_2\subseteq\cdots,\qquad A=\bigcup_{N\ge 0}A_N,$$
such that
\begin{enumerate}
\item each $A_N$ is a finitely generated free $R$-module and each quotient
      $A_N/A_{N-1}$ is a free $R$-module;
\item $A_M A_N\subseteq A_{M+N}$ for all $M,N\ge 0$; and
\item $A$ is generated as an $R$-algebra by finitely many elements
      $g_1,\dots,g_r$, all lying in $A_d$ for some fixed $d$.
\end{enumerate}
If $A\otimes_R F$ has  a central element not contained in $F\cdot 1$, then $A\otimes_R K$ has  a central element not contained in $K\cdot 1$.
\end{lemma}

\begin{proof}
Since the $g_i$ generate $A$ as an $R$-algebra, an element $a\in A$ is central if and only if
$[a,g_i]=0$ for $i=1,\dots,r$; the same condition holds for $A\otimes_R K$, since the
images $g_i\otimes 1$ generate $A\otimes_R K$ as a $K$-algebra.  Now suppose $a\in A_N\setminus F\cdot 1$ is central.  For $b\in A_N$ we have
$[b,g_i]\in A_{N+d}$ by assumption $(2)$, so there is an $R$-linear map of
finitely generated free $R$-modules
$\Phi\colon A_N\longrightarrow (A_{N+d})^{r}$ given by 
$b\mapsto\bigl([b,g_1],\dots,[b,g_r]\bigr)$.
By construction
$Z(A)\cap A_N=\ker\Phi$. 

Fix $R$-bases of $A_N$ and of $(A_{N+d})^{r}$, and let $q$ and $p$ denote the sizes of the respective bases. We can then represent $\Phi$ as a $p\times q$ matrix $Y$ with entries in $R$, and we regard it as a matrix over $F$.  Let
$\rho:=\operatorname{rank}_F(Y)$.  Since $1,a\in {\rm ker}(\Phi)$ and they are linearly independent over $F$, $\rho\le q-2$.  Hence every $(q-1)\times(q-1)$ minor
of $Y$ vanishes. If we let $\theta(Y)\in M_{p\times q}(K)$ denote the matrix obtained by applying $\theta$ entry-wise, we then see that every $(q-1)\times (q-1)$ minor of $\theta(Y)$ vanishes and so $\theta(Y)$ has rank at most $q-2$ as a matrix over $K$.
Consequently
$$\dim_K\ker\!\bigl(\Phi\otimes {\rm Id}_K)
   \ge 2$$
   and so we see that $A\otimes_R K$ necessarily has a non-scalar central element in $A_N\otimes_R K$. 

\end{proof}



\begin{lemma}\label{lem:34}
Let $K$ be a field and let $A$ be an algebra from either family (3) or (4) from Theorem \ref{thm:main}.  Then there exists a central element of $A$ that is not in $K$.
\end{lemma}
\begin{proof}
For the algebras from family (3), we write $u(x)=c_0+c_1x+c_2x^2+c_3x^3+c_4x^4$.  Using Sage Math,\footnote{The element was found by solving for a lift of Poisson central elements for the corresponding Poisson algebras.} we always have a central element
$C:=C_6+C_5+C_4+C_3+C_2+C_1$ with
$$C_6=20\delta y^3-30z^2, \quad C_5= -12c_4x^5-20Ax^3y+30Ax^2z-20\delta Axy^2+10\delta A yz,$$
\begin{align*} C_4 &=-\big(8\delta^2c_4^2-3\beta\delta c_4+15c_3\big)x^4-10\big(18\beta^2\delta-11\beta\delta^2c_4+3\delta c_3-6\gamma_2\big)x^2y\\
&+30\big(5\beta^2\delta-3\beta\delta^2c_4+\delta c_3-2\gamma_2\big)xz-10\delta\big(4\beta^2\delta-\beta\delta^2c_4+\delta c_3-4\gamma_2\big)y^2,\end{align*}
$$C_3=-\big(8\delta^2c_3c_4-21\beta^2\delta^2c_4-35\beta\delta c_3+56\delta\gamma_2c_4+20c_2\big)x^3-20Bxy+10Bz,$$
\begin{align*} C_2 &=-\big(8\delta^2c_2c_4-3\beta^3\delta^3c_4-5\beta^2\delta^2c_3-7\beta\delta^2\gamma_2c_4-40\beta\delta c_2\\
&+54\delta\gamma_1c_4+15\delta\gamma_2c_3+30c_1\big)x^2 +10\big(\delta^2\gamma_1c_4-\delta c_1+6\gamma_0\big)y\end{align*}
$$C_1=-\big(8\delta^2c_1c_4-9\beta\delta^2\gamma_1c_4-30\beta\delta c_1+48\delta\gamma_0c_4+15\delta\gamma_1c_3+60c_0\big)x$$

where $A:=2\delta c_4-3\beta$ and $$B:=-3\beta^3\delta^2+2\beta^2\delta^3c_4-\beta\delta^2c_3+3\beta\delta\gamma_2-\delta^2\gamma_2c_4+\delta c_2-3\gamma_1.$$


For the algebras from family (4), we have relations
$$[x,y]=z, [x,z]= z -2 xy +\zeta, [y,z]= y^2 +u(x),$$
with $u(x)\in K[x]$.  We let $n$ denote the degree of $u(x)$.  We then create commuting indeterminates $t_0,\ldots ,t_n, s$. We let
$R=\mathbb{Z}[t_0,\ldots ,t_n,s]$, $F={\rm Frac}(R)$, and let $\theta: R\to K$ be the ring homomorphism given by sending $s$ to $\zeta$ and $t_i$ to the coefficient of $x^i$ in $u(x)$. We now consider the $R$-algebra $A$ with generators $x,y,z$ and relations
$[x,y]=z, [x,z]=z-2xy+s, [y,z] = y^2+t_0 + \cdots +t_n x^n$, then an application of Bergman's diamond lemma \cite{Berg}, using the fact that the Jacobi identity holds, gives that $A$ is a free $R$-algebra with basis $x^i y^j z^k$.  Then by Lemma \ref{lem:lift} it suffices to show that $A\otimes_R F$ has a non-trivial central element.

So we now work with the algebra $A\otimes_R F$ and the field $F$.  We note that ${\rm ad}_x$ preserves the left $F[x]$-submodule of our algebra $F[x]+F[x]z+F[x]y$.  From here we can produce a non-trivial element $$\tilde{C}:=z^2+ 2xy^2 - 2yz+  2y^2 -2s y$$ that commutes with $x$. We found this element using Sage Math.

Now we claim there is a polynomial $f(x)\in F[x]$ such that $\tilde{C}+f(x)$ is central.  Since our algebras are generated by $x$ and $y$ and since $\tilde{C}+f(x)$ commutes with $x$, it suffices to find $f(x)$ such that $[\tilde{C}+f(x),y]=0$.

Then \begin{align*}
[\tilde{C},y] &= -(y^2 + u(x)) z - z(y^2+u(x)) +2zy^2 + 2y(y^2+u(x)) \\
&=-y^2z+zy^2 -u(x)z - zu(x) +2y^3 +2y u(x).
\end{align*}
Since $y^2z-zy^2=[y^2,z] = [y,z] y + y[y,z] = 2y^3 + y u(x)+u(x)y$, we see that 
$[\tilde{C},y] = yu(x)-u(x)y -u(x)z - zu(x) = [-u(x),y] - u(x) z - zu(x)$. 
 Thus it suffices to show that $u(x)z+zu(x) = [g(x),y]$ for some polynomial $g(x)$.  Furthermore, since $z=[x,y]$ we can write $u(x)z+zu(x)$ as
 $u(x) xy - u(x)yx + xy u(x) - y x u(x) = [u(x)x ,y] + x yu(x)-  u(x)yx$.
 Thus it suffices to show there is some polynomial $h(x)$ such that $[h(x),y] = x y u(x)- u(x)y x$.  By linearity, it suffices to do this when $u(x)=x^n$ with $n\ge 0$ and so we must show that for each $n\ge 0$ there is a polynomial $q_n(x)$ such that $[q_n(x),y] = xyx^n- x^n yx$.


To do this, observe that our algebra is a left $F[s,t]$-module via the action $s\cdot a = x \cdot a, t\cdot a= a\cdot x$ for $a$ in our algebra.  Then let $I$ be the ideal of $F[s,t]$ consisting of all elements that annihilate $y$.  Since 
$[x,[x,y]] = z- 2xy+\zeta = [x,y]- 2xy+\zeta$, we see that $(s-t)((s-t)^2 -(s-t) + 2s)\in I$.  Then $I$ contains the ideal generated by $(s-t)((s-t)^2+s+t)$.
We let $J$ denote the ideal generated by $(s-t)^2+s+t$.

Our goal is now to show that for each $n\ge 0$ there is a polynomial $q_n$ such that
$q_n(s)-q_n(t) - st^n + s^nt\in I$.  Since $q_n(s)-q_n(t)-st^n+s^n t$ is necessarily divisible by $(s-t)$ and since $(s-t)$ and $(s-t)^2+s+t$ are coprime, it suffices to show that for each $n\ge 0$ there is a polynomial $q_n$ such that
$q_n(s)-q_n(t) - st^n + s^nt\in J$. 

Now in characteristic zero, the curve $(s-t)^2+s+t=0$ is irreducible of genus $0$ and can be parametrized by $s=(T-T^2)/2, t=(-T-T^2)/2$, with $T=s-t$ and we have a $F$-algebra isomorphism $\Psi: F[s,t]/J\to F[T]$ given by the parametrization just given. Now $F[T]$ has a $F$-algebra automorphism $\sigma$ of order $2$ given by $T\mapsto -T$, and this automorphism induces an automorphism $\tilde{\sigma}$ of $F[s,t]/J$ that interchanges the images of $s$ and $t$.  Then we claim that 
$S=\Psi(\{q(s)-q(t) \colon q(\theta)\in F[\theta]\})$ is precisely the set of odd polynomials in $T$.  To see this, observe that $q(s)-q(t)$ is sent to its negative under $\tilde{\sigma}$ and so every element of $S$ is sent to its negative under $\sigma$ and hence is odd. On the other hand, $\Psi(s^n-t^n)$ is a degree $2n-1$ polynomial with leading coefficient $2^{-n+1}n$ and so $S$ contains a basis for the odd polynomials in $T$.  Thus to obtain the result, it suffices to note that $st^n-s^n t = st (t^{n-1}-s^{n-1})$ is sent to its negative by $\tilde{\sigma}$ and hence has odd image in $F[T]$. Thus we have proven the existence of $f(x)$ such that $[\tilde{C}+f(x),y]=[\tilde{C}+f(x),x]=0$ and by construction we in fact have $f(x)\in \mathbb{Q}[t_0,\ldots ,t_n][x]$.  Thus we have proved the existence of a non-trivial central element for algebras in family (4).
\end{proof}
\begin{proof}[Proof of Theorem \ref{thm:main2}]

The families of type (1) are iterated Ore extensions.  A unipotent automorphism has order a power of $p$ when the base field has characteristic $p>0$ and so the fact that these algebras are PI follows from work of Brown and Zhang \cite[Theorem 3.3]{BZ}. 

For the remaining families, in light of Proposition \ref{prop:invoke}, it suffices to show that each family from (2)--(9) has a non-trivial centre in characteristic $p>0$.
For the algebras from family (2) with generators $x,y,z$ and relations $[x,y]=z, \qquad [x,z]=\gamma y, \qquad [y,z]=f(x)$ for some $\gamma\in K$ and some $f(x)\in K[x]$, we have
${\rm ad}_x^n(y) = \gamma^{(n-1)/2} z$, and ${\rm ad}_x^n(z)=\gamma^{(n+1)/2}y$ if $n$ is odd; and ${\rm ad}_x^n(y)=\gamma^{n/2} y, {\rm ad}_x^n(z)=\gamma^{n/2} z$ if $n$ is even.
In particular, since ${\rm ad}_x^{p^i} = {\rm ad}_{x^{p^i}}$ we see that for all fields of characteristic $p>0$, $x^{p^2}-\gamma^{(p^2-p)/2} x^p$ is central.  Thus we obtain the desired result for the algebras in (2).


Families (3) and (4) have a non-trivial centre by Lemma \ref{lem:34} and thus these families satisfy polynomial identities in positive characteristic.

Families (5), (7), and (9) have relations of the form $[x,y]=z$ and either $[y,z]=c z$ or $[x,z]=cz$ and so either ${\rm ad}_y^n(x) = -c^{n-1} z$ for all $n\ge 1$ or ${\rm ad}_x^n(y)=c^{n-1}z$ for $n\ge 1$, and so in these algebras we have either $y^p - c^{p-1} y$ or $x^p-c^{p-1}x$ is central. 

For family (6) we have characteristic $2$. When $\alpha=0$ we have relations
$$[x,y]=z, \quad [x,z]= \beta y +  \gamma_2 x^2 + \gamma_1 x + \gamma_0 +\delta y^2,\quad [y,z]=u(x)+\gamma_1y +\gamma_2 z,$$
It follows that ${\rm ad}_x$ preserves the $K$-vector space spanned by $1, x, x^2, x^3, x^4, y, z, y^2$.  Since the space of $K$-linear endomorphisms of $V$ is $64$-dimensional, since ${\rm ad}_x^{p^i}={\rm ad}_{x^{p^i}}$ for $i\ge 1$, it follows that there is a non-trivial $K$-linear dependence between 
${\rm ad}_x, {\rm ad}_{x^p}, \ldots , {\rm ad}_{x^{p^{64}}}$ when we restrict the action to $V$.  Consequently there is a non-trivial polynomial in $x$ that commutes with $y$ and $z$ and hence is central. When $\alpha=1$, we have the relations
$$[x,y]=z, [x,z]=z + \beta y + (\delta^{-1} x^3+ \gamma_2 x^2 + \gamma_1 x + \gamma_0 ) +\delta y^2,$$
and
$$[y,z] = u(x) + (\delta^{-1} x^2+\gamma_1+\beta \delta^{-1})y + (\delta^{-1} x+\gamma_2+\delta^{-1})z.$$
Then we see that the $K(x)$-vector subspace of the quotient division algebra of our ring that is spanned by $1, y, z, y^2$ is invariant under ${\rm ad}_x$ and moreover, we have
$${\rm ad}_x(1)=0, {\rm ad}_x(y)=z, {\rm ad}_x(z)= z + \beta y + (\delta^{-1} x^3+ \gamma_2 x^2 + \gamma_1 x + \gamma_0 ) +\delta y^2$$
and $${\rm ad}_x(y^2) = {\rm ad}_y(z)=u(x) + (\delta^{-1} x^2+\gamma_1+\beta \delta^{-1})y + (\delta^{-1} x+\gamma_2+\delta^{-1})z.$$
Then if we use this fact we can verify that $$h(x)=x^4 + x^3 + (1+\beta+\delta\gamma_2)x^2 + (\beta+\delta\gamma_1)x$$ is central and so we are done.

Finally for family (8), ${\rm ad}_x$ preserves the $K$ vector space spanned by $1,x,x^2,x^3,y,z$ and so we obtain a centre using the argument employed in handling case (6), but now using invariance of a $6$-dimensional space to produce a central polynomial in $x$.
\end{proof}
\begin{thebibliography}{99}
\bibitem{ASZ} M. Artin, L. W. Small, and J. J. Zhang, 
Generic flatness for strongly Noetherian algebras.
\emph{J. Algebra} {\bf 221} (1999), no. 2, 579–610.
\bibitem{Belldiv} J. P. Bell, Transcendence degree of division algebras. 
\emph{Israel J. Math.} {\bf 190} (2012), 195–211.

\bibitem{Bell} J. P. Bell, Filtered deformations of commutative algebras of Krull dimension two.
\emph{Math. Z.} {\bf 307} (2024), no. 2, Paper No. 27, 9 pp.
\bibitem{Berg} G. M. Bergman, 
The diamond lemma for ring theory.
\emph{Adv. in Math.} {\bf 29} (1978), no. 2, 178–218.

\bibitem{BZ} K. A. Brown and J. J. Zhang, Iterated Hopf Ore extensions in positive characteristic.
\emph{J. Noncommut. Geom.} {\bf 16} (2022), no. 3, 787–837.
\bibitem{CEW}  J. Cuadra, P. Etingof, and C. Walton, Semisimple Hopf actions on Weyl
algebras. \emph{Adv. Math.} {\bf 282} (2015), 47--55.
%%(QUESTION on Page 2 asks if gr(B) mod $p$ is PI if B is a finitely geenerated domain in char $p$)Let B be a Z+-filtered algebra over k with gr(B) a finitely generated commutative domain. Does (†) hold for any large prime p?


\bibitem{Etingof} P. Etingof, 
A PI degree theorem for quantum deformations. \emph{J. Algebra} {\bf 466} (2016), 308–313.
\bibitem{Lam} T. Y. Lam, 
\emph{A first course in noncommutative rings.} Second edition. Grad. Texts in Math., 131
Springer-Verlag, New York, 2001.

\bibitem{Li} B. Li, \emph{Filtered deformations of trivariate
polynomial rings}, Master's Project, Pure Mathematics Department, University of Waterloo, August 2026. Available online at: {\tt https://www.cadenze.ca/files/pub/mmath-thesis.pdf}.
\bibitem{Vam} P. V\'amos, On the minimal prime ideals of a tensor product of two fields. \emph{Math. Proc. Cambridge Philos. Soc.} {\bf 84} (1978), no. 1, 25--35.

%\bibitem{ZSvol1} 
%O. Zariski and P. Samuel, 
%\emph{Commutative algebra. Vol. 1.
%With the cooperation of I. S. Cohen.} Corrected reprinting of the 1958 edition
%Grad. Texts in Math., No. 28
%Springer-Verlag, New York-Heidelberg-Berlin, 1975.

\end{thebibliography}


%\begin{thebibliography}(99}
%\bibitem{ASZ} Artin, M.(1-MIT); Small, L. W.(1-UCSD); Zhang, J. J.(1-WA)
%Generic flatness for strongly Noetherian algebras.
%J. Algebra 221 (1999), no. 2, 579–610.
%\bibitem{Vam} P. V\'amos, On the minimal prime ideal of a tensor product of two fields.Math. \emph{Proc. Cambridge Philos. Soc.} {\bf 84} (1978), no. 1, 25--35.

%\end{thebibliography}
 \end{document}

%MSC codes.Maybe add a polynomial-identity-rings classification (16R10? 16R20?) to the subject classification.

